\documentclass[10pt,a4paper]{amsart}
\usepackage[margin=19mm]{geometry}
\usepackage{amsmath,amssymb,mathtools}
\usepackage[scr=rsfs,cal=euler]{mathalfa}
\usepackage{xfrac}
\usepackage[unicode,hidelinks,bookmarks=false]{hyperref}
\newcommand{\ie}{i.e.\ }
\newcommand{\cf}{cf.\ }
\newcommand{\resp}{respectively\ }
\newcommand{\Subsection}{\subsection}
\providecommand{\nobreakdash}{\nobreak\hbox{-}}
\setlength{\emergencystretch}{2em}
\multlinegap0pt
\usepackage{amssymb}
\usepackage[shortlabels]{enumitem}
\newtheorem{lemma}{Lemma}
\newtheorem{proposition}{Proposition}
\newcommand{\D}{\mathrm{D}}
\newcommand{\W}{\mathrm{W}}
\newcommand{\id}{{\mathrm{id}}}
\title{On deformations of foliations and characteristic classes}
\author{Taro Asuke}
\date{Source: arXiv:2603.23827v1 (CC BY 4.0)}
\begin{document}
\maketitle

\noindent\emph{Source and licence.} On deformations of foliations and characteristic classes. Version arXiv:2603.23827v1 (CC BY 4.0), distributed by arXiv under the Creative Commons Attribution 4.0 International licence, http://creativecommons.org/licenses/by/4.0/, as recorded in the arXiv licence metadata for this exact version and shown on the abstract page https://arxiv.org/abs/2603.23827v1. Full licence notice: \texttt{licenses/arxiv-2603.23827-CC-BY-4.0.txt}.\par\smallskip

\noindent\textit{Editorial scope.} Counted unit: the article's proposition prop10.6 (source0.tex lines 1937--1940). The article's complete original proof of it is delivered with it (source0.tex lines 1941--2061, one \texttt{proof} environment of 121 lines). No mathematics is written, repaired or re-derived: the statement and the proof are the article's own lines, in the article's own notation, and no retained line is substituted or re-flowed.  Retained uncounted as the source-local context the proof needs: lines 1917--1920.  Nothing was omitted from any retained span, so the package carries no omission record.  No retained line contains a citation, so the wrapper carries no bibliography and no bibliography entry.  No retained line contains a figure environment, an image-inclusion command or drawing code, so no image is delivered and the package declares no asset dependency.  Editorial formatting only, in addition: an AMS \texttt{amsart} preamble that restates the article's own declarations and macros used by the retained lines, namely the environments \texttt{lemma}, \texttt{proposition} on separate counters, together with the standard packages those lines need.  The retained environments print in this document as Lemma 1, Proposition 1 in order of appearance, whereas the article prints the same items with the numbering of its own full document.  The complete native source of the article as distributed in the arXiv e-print for this exact version is bundled unchanged as \texttt{sources/197111/source0.tex}.
\begin{lemma}
\label{lem10.8}
Elements of the form $c_{(i)}c_{(j)}c_{(k)}$ with $i+j+k\leq5$ are trivial in~$\D^r\W_1$ if $r\geq5$.
\end{lemma}

\begin{proposition}
\label{prop10.6}
We have $F_{0,k}\subset H^*(\D^\infty\W_1)\cong\{0\}$\, for $k=2,3,4$.
\end{proposition}

\begin{proof}
For a while, $c,\dot{c},\ddot{c},\dddot{c}$ denote $c_{1,(0)},c_{1,(1)},c_{1,(2)},c_{1,(3)}$, respectively.
We make use a similar notations for $h$.
Recall that $p_{1,k}$ denotes the projection from $H^*(\D^\infty\W_1)^k$ to~$F_{0,k}$.

Let $k=2$.
As $c^2=c\dot{c}=0$, the following elements gives rise to a set of basis for $(\D^\infty\W_1)^2$:
\[
\ddot{h}, h\ddot{h}, \ddot{h}c, h\ddot{h}c,%
\ddot{c}, \ddot{c}h, \ddot{c}c, \ddot{c}hc,%
\dot{h}\dot{c}, h\dot{h}\dot{c},%
\dot{c}^2, \dot{c}^2h.
\]
The images of these elements by $p_{1,2}$ to the $0$-eigenspace of $\delta\sigma$ are equal~to
\[
0, 0, -\dot{h}\dot{c},-h\dot{h}\dot{c},%
0, -\dot{c}\dot{h}, \ddot{c}c, \ddot{c}hc,%
\dot{h}\dot{c},h\dot{h}\dot{c},%
\dot{c}^2,\dot{c}^2h.
\]
Hence cocycles in the $0$-eigenspace $F_{0,2}$ of $\delta\sigma$ are linear combinations of
\[
\dot{c}^2,\ddot{c}c,\dot{c}^2h,\ddot{c}hc.
\]
As we have $\dot{c}^2h=d(\dot{c}\dot{h}h)$ and $\ddot{c}hc=\ddot{h}hc$, all of these elements are exact.

Next, we study $H^*(\D^\infty\W_1)^3$.
We have
\begin{align*}
&p_{1,3}(\dddot{h})=0,\\*
&p_{1,3}(\dddot{h}c)=-\frac12\ddot{h}\dot{c}+\frac12\dot{h}\ddot{c}=-\frac12d(\dot{h}\ddot{h}),\\*
&p_{1,3}(h\dddot{h}c)=-\dot{h}\ddot{h}c-\frac12h\ddot{h}\dot{c}+\frac12h\dot{h}\ddot{c}=-\frac32\dot{h}\ddot{h}c+\frac12d(h\dot{h}\ddot{h}),\\*
&p_{1,3}(\dddot{c})=0,\\*
&p_{1,3}(h\dddot{c})=-\frac12\dot{h}\ddot{c}+\frac12\ddot{h}\dot{c}=\frac12d(\dot{h}\ddot{h}),\\*
&p_{1,3}(h\dddot{c}c)=-\frac12\dot{h}\ddot{c}c+\dot{h}\dot{c}^2+\frac12\ddot{h}\dot{c}c=\dot{h}\dot{c}^2+\frac12d(\dot{h}\ddot{h}c),\\
% \end{align*}
% Similarly, we have
% \begin{align*}
&p_{1,3}(\dot{h}\ddot{h})=\dot{h}\ddot{h},\\*
&p_{1,3}(\dot{h}\ddot{h}c)=\dot{h}\ddot{h}c,\\*
&p_{1,3}(h\dot{h}\ddot{h})=h\dot{h}\ddot{h},\\*
&p_{1,3}(h\dot{h}\ddot{h}c)=h\dot{h}\ddot{h}c,\\*
&p_{1,3}(\dot{h}\ddot{c})=-\frac12\ddot{h}\dot{c}+\frac12\dot{h}\ddot{c}=-\frac12d(\dot{h}\ddot{h}),\\*
&p_{1,3}(\dot{h}\ddot{c}c)=\dot{h}\ddot{c}c=-\dot{h}\dot{c}^2+\dot{h}(\ddot{c}c+\dot{c}^2)=-\dot{h}\dot{c}^2\\*
&p_{1,3}(\dot{h}\dot{c}^2)=\dot{h}\dot{c}^2,\\*
&p_{1,3}(h\ddot{c}\dot{c})=-\frac12\dot{h}\dot{c}^2,\\
&p_{1,3}(h\dot{c}^3)=h\dot{c}^3=-h\ddot{c}c\dot{c}+h(\ddot{c}c+\dot{c}^2)\dot{c}=0.
\end{align*}
Hence $F_{0,3}$ is generated by $\dot{h}\dot{c}^2$.
On the other hand, we have $\dot{h}\dot{c}^2=-\dot{h}\ddot{c}c=d(\dot{h}\ddot{h}c)$.
Therefore, we have $H^*(\D^\infty\W_1)^3\cong\{0\}$.

Finally, we study $H^*(\D^\infty\W_1)^4$.
% In what follows, $c$ denotes $c_1$ and $c^{(i)}$ denotes $c_{1,(i)}$ for $i\geq1$.
% We make use of similar notations for $h$ and $h^{(i)}$.
We have
\[
p_{1,4}=\id-\frac13\delta\circ\sigma+\frac1{15}\delta^2\circ\sigma^2-\frac1{90}\delta^3\circ\sigma^3,
\]
which is denoted by $p$ for simplicity.

We have
\begin{align*}
&p(h_{(4)})=0,\\*
&p(h_{(4)}c_{(0)})=-\frac15h_{(3)}c_{(1)}+\frac35h_{(2)}c_{(2)}-\frac15h_{(1)}c_{(3)},\\*
&p(h_{(0)}h_{(4)})=0,\\*
&p(h_{(0)}h_{(4)}c_{(0)})=-\frac15h_{(0)}h_{(3)}c_{(1)}+\frac95h_{(1)}h_{(2)}c_{(1)}+\frac35h_{(0)}h_{(2)}c_{(2)}-\frac15h_{(0)}h_{(1)}c_{(3)},\\*
&p(h_{(1)}h_{(3)})=0,\\*
&p(h_{(1)}h_{(3)}c_{(0)})=-h_{(1)}h_{(2)}c_{(1)},\\*
&p(h_{(0)}h_{(1)}h_{(3)})=0,\\*
&p(h_{(0)}h_{(1)}h_{(3)}c_{(0)})=-h_{(0)}h_{(1)}h_{(2)}c_{(1)},\\
&p(h_{(3)}c_{(1)})=\frac15h_{(3)}c_{(1)}-\frac35h_{(2)}c_{(2)}+\frac15h_{(1)}c_{(3)},\\*
&p(h_{(0)}h_{(3)}c_{(1)})=\frac15h_{(0)}h_{(3)}c_{(1)}-\frac45h_{(1)}h_{(2)}c_{(1)}-\frac35h_{(0)}h_{(2)}c_{(2)}+\frac15h_{(0)}h_{(1)}c_{(3)},\\*
&p(h_{(1)}h_{(2)}c_{(1)})=h_{(1)}h_{(2)}c_{(1)},\\*
&p(h_{(0)}h_{(1)}h_{(2)}c_{(1)})=h_{(0)}h_{(1)}h_{(2)}c_{(1)},\\
&p(h_{(2)}c_{(2)})=-\frac15h_{(3)}c_{(1)}+\frac35h_{(2)}c_{(2)}-\frac15h_{(1)}c_{(3)},\\*
&p(h_{(2)}c_{(2)}c_{(0)})=-\frac15h_{(3)}c_{(0)}c_{(1)}+\frac35h_{(2)}c_{(0)}c_{(2)}-\frac1{15}h_{(2)}c_{(1)}{}^2-\frac15h_{(1)}c_{(0)}c_{(3)}+\frac1{15}h_{(1)}c_{(1)}c_{(2)};\\*
&\quad\text{note that we have}\\*
&\quad h_{(2)}c_{(2)}c_{(0)}\begin{aligned}[t]
&=-h_{(2)}c_{(1)}{}^2\\*
&=d(h_{(1)}h_{(2)}c_{(1)})+h_{(1)}c_{(1)}c_{(2)}\\*
&=d(h_{(1)}h_{(2)}c_{(1)})-\frac13h_{(1)}c_{(0)}c_{(3)}\\*
&=d(h_{(1)}h_{(2)}c_{(1)})+\frac13d(h_{(1)}h_{(3)}c_{(0)}),
\end{aligned}\\*
&p(h_{(2)}c_{(1)}{}^2)=\frac23h_{(2)}c_{(1)}{}^2-\frac23h_{(1)}c_{(1)}c_{(2)},\\*
&p(h_{(0)}h_{(2)}c_{(2)})=-\frac15h_{(0)}h_{(3)}c_{(1)}+\frac35h_{(0)}h_{(2)}c_{(2)}-\frac15h_{(1)}h_{(2)}c_{(1)}-\frac15h_{(0)}h_{(1)}c_{(3)},\\*
&p(h_{(0)}h_{(2)}c_{(0)}c_{(2)})=\frac13h_{(0)}h_{(2)}c_{(0)}c_{(2)}-\frac13h_{(0)}h_{(2)}c_{(1)}{}^2-\frac13h_{(0)}h_{(1)}c_{(3)}c-\frac13h_{(0)}h_{(1)}c_{(1)}c_{(2)},\\*
&p(h_{(0)}h_{(2)}c_{(1)}{}^2)=\frac23h_{(0)}h_{(2)}c_{(1)}{}^2-\frac23h_{(0)}h_{(1)}c_{(1)}c_{(2)},\\
&p(h_{(1)}c_{(2)}c_{(1)})=-\frac13h_{(2)}c_{(1)}{}^2+\frac13h_{(1)}c_{(1)}c_{(2)},\\*
&p(h_{(0)}h_{(1)}c_{(1)}c_{(2)})=-\frac13h_{(0)}h_{(2)}c_{(1)}{}^2+\frac13h_{(0)}h_{(1)}c_{(1)}c_{(2)},\\
&p(h_{(0)}c_{(2)}{}^2)=\frac2{15}h_{(2)}c_{(1)}{}^2-\frac2{15}h_{(1)}c_{(1)}c_{(2)}+\frac35h_{(0)}c_{(2)}{}^2-\frac25h_{(0)}c_{(1)}c_{(3)};\\*
&\quad\text{note that we have}\\*
&\quad
h_{(0)}c_{(2)}{}^2=-h_{(0)}c_{(1)}c_{(1)}=d(h_{(0)}h_{(1)}c_{(1)}),\\*
&p(h_{(0)}c_{(1)}{}^2c_{(2)})=-\frac13h_{(1)}c_{(1)}{}^3=0,\\*
&p(h_{(0)}c_{(0)}c_{(2)}{}^2)=h_{(0)}c_{(0)}c_{(2)}{}^2=0,
\end{align*}
where we make use of Lemma~\ref{lem10.8}.

Therefore, cochains in $E_{0,4}$ are prepresented by linear combinations of the following elements:
\begin{align*}
&h_{(3)}c_{(1)}-3h_{(2)}c_{(2)}+h_{(1)}c_{(3)},\\*
&h_{(0)}h_{(3)}c_{(1)}-3h_{(0)}h_{(2)}c_{(2)}+h_{(0)}h_{(1)}c_{(3)},\\*
&h_{(1)}h_{(2)}c_{(1)},\\*
&h_{(0)}h_{(1)}h_{(2)}c_{(1)},\\*
&{}-3h_{(2)}c_{(0)}c_{(2)}+h_{(1)}c_{(0)}c_{(3)},\quad\text{(exact)}\\*
&h_{(2)}c_{(1)}{}^2-h_{(1)}c_{(1)}c_{(2)},\quad\text{(exact)}\\*
&-3h_{(0)}h_{(2)}c_{(0)}c_{(2)}+h_{(0)}h_{(1)}c_{(0)}c_{(3)},\quad\text{(exact)}\\*
&h_{(0)}h_{(2)}c_{(1)}{}^2-h_{(0)}h_{(1)}c_{(1)}c_{(2)},\quad\text{(exact)}\\*
&3h_{(0)}c_{(2)}{}^2-2h_{(0)}c_{(1)}c_{(3)}.\quad\text{(exact)}
\end{align*}
We can show that the above elements except the first four are exact.
Hence, only the second and the third one of the above elements can yield non-trivial classes.
The derivative of them are equal to
\begin{align*}
&h_{(0)}c_{(1)}c_{(3)}-3h_{(0)}c_{(2)}{}^2+3h_{(2)}c_{(0)}c_{(2)}+h_{(0)}c_{(1)}c_{(3)}-h_{(1)}c_{(0)}c_{(3)}\\*
&=2h_{(0)}c_{(1)}c_{(3)}+3h_{(2)}c_{(0)}c_{(2)}-h_{(1)}c_{(0)}c_{(3)},\\*
&h_{(1)}c_{(1)}c_{(2)}-h_{(2)}c_{(1)}{}^2.
\end{align*}
It is easy to see that they are linearly independent so that we have $F_{0,4}\cong\{0\}$.
\end{proof}

\end{document}
